% Options for packages loaded elsewhere
\PassOptionsToPackage{unicode}{hyperref}
\PassOptionsToPackage{hyphens}{url}
\documentclass[
  ignorenonframetext,
]{beamer}
\newif\ifbibliography
\usepackage{pgfpages}
\setbeamertemplate{caption}[numbered]
\setbeamertemplate{caption label separator}{: }
\setbeamercolor{caption name}{fg=normal text.fg}
\beamertemplatenavigationsymbolsempty
% remove section numbering
\setbeamertemplate{part page}{
  \centering
  \begin{beamercolorbox}[sep=16pt,center]{part title}
    \usebeamerfont{part title}\insertpart\par
  \end{beamercolorbox}
}
\setbeamertemplate{section page}{
  \centering
  \begin{beamercolorbox}[sep=12pt,center]{section title}
    \usebeamerfont{section title}\insertsection\par
  \end{beamercolorbox}
}
\setbeamertemplate{subsection page}{
  \centering
  \begin{beamercolorbox}[sep=8pt,center]{subsection title}
    \usebeamerfont{subsection title}\insertsubsection\par
  \end{beamercolorbox}
}
% Prevent slide breaks in the middle of a paragraph
\widowpenalties 1 10000
\raggedbottom
\AtBeginPart{
  \frame{\partpage}
}
\AtBeginSection{
  \ifbibliography
  \else
    \frame{\sectionpage}
  \fi
}
\AtBeginSubsection{
  \frame{\subsectionpage}
}
\usepackage{iftex}
\ifPDFTeX
  \usepackage[T1]{fontenc}
  \usepackage[utf8]{inputenc}
  \usepackage{textcomp} % provide euro and other symbols
\else % if luatex or xetex
  \usepackage{unicode-math} % this also loads fontspec
  \defaultfontfeatures{Scale=MatchLowercase}
  \defaultfontfeatures[\rmfamily]{Ligatures=TeX,Scale=1}
\fi
\usepackage{lmodern}
\ifPDFTeX\else
  % xetex/luatex font selection
\fi
% Use upquote if available, for straight quotes in verbatim environments
\IfFileExists{upquote.sty}{\usepackage{upquote}}{}
\IfFileExists{microtype.sty}{% use microtype if available
  \usepackage[]{microtype}
  \UseMicrotypeSet[protrusion]{basicmath} % disable protrusion for tt fonts
}{}
\makeatletter
\@ifundefined{KOMAClassName}{% if non-KOMA class
  \IfFileExists{parskip.sty}{%
    \usepackage{parskip}
  }{% else
    \setlength{\parindent}{0pt}
    \setlength{\parskip}{6pt plus 2pt minus 1pt}}
}{% if KOMA class
  \KOMAoptions{parskip=half}}
\makeatother
\usepackage{longtable,booktabs,array}
\newcounter{none} % for unnumbered tables
\usepackage{calc} % for calculating minipage widths
\usepackage{caption}
% Make caption package work with longtable
\makeatletter
\def\fnum@table{\tablename~\thetable}
\makeatother
\setlength{\emergencystretch}{3em} % prevent overfull lines
\providecommand{\tightlist}{%
  \setlength{\itemsep}{0pt}\setlength{\parskip}{0pt}}
\usepackage{bookmark}
\IfFileExists{xurl.sty}{\usepackage{xurl}}{} % add URL line breaks if available
\urlstyle{same}
\hypersetup{
  hidelinks,
  pdfcreator={LaTeX via pandoc}}

\author{\texorpdfstring{}{}}
\date{}

\begin{document}

\begin{frame}[fragile]{Tooling and Infrastructure: Software Ecosystem,
Knowledge Graph Deployment, and Quality Assurance \label{sec:tooling}}
\protect\phantomsection\label{tooling-and-infrastructure-software-ecosystem-knowledge-graph-deployment-and-quality-assurance}
The practical utility of a computational meta-analysis depends on robust
tooling at each pipeline stage: assertion extraction, modeling and
simulation, knowledge graph infrastructure, and quality assurance.

\begin{block}{LLM-Based Assertion Extraction}
\protect\phantomsection\label{llm-based-assertion-extraction}
Extracting structured assertions from unstructured text is the most
labor-intensive component of knowledge graph construction. Manual
annotation produces high-quality results but does not scale to corpora
of thousands of papers---a constraint demonstrated by Knight et
al.~\citep{knight2022fep}, whose systematic literature analysis of FEP
and Active Inference publications required manual coding of structural,
visual, and mathematical features for hundreds of annotated papers. We
implement a hybrid approach: LLMs perform initial extraction, with human
review for validation and correction.

Our extraction pipeline deploys a locally hosted LLM through Ollama
\citep{ollama2024}. Each paper's abstract is assessed against the eight
hypothesis definitions in a structured prompt requesting a JSON array of
assessments. Unlike keyword matching, which detects only topical terms,
the LLM evaluates the \emph{semantic relationship} between a paper's
claims and each hypothesis. Papers critiquing the FEP correctly receive
``contradicts'' assessments for FEP Universality (H1), while methodology
tutorials receive ``neutral'' assessments reflecting their pedagogical
character. Detailed prompt engineering, schemas, and failure modes are
documented in the supplementary extraction pipeline (see
\hyperref[sec:extraction_pipeline]{Section 3}).
\end{block}

\begin{block}{Software Ecosystem}
\protect\phantomsection\label{software-ecosystem}
The Active Inference community has developed several specialized
software tools, though the ecosystem remains highly fragmented---no
single package spans the full spectrum from theoretical simulation to
empirical data analysis:

\textbf{pymdp.} The pymdp library \citep{heins2022pymdp} provides a
Python implementation of active inference for discrete state-space
POMDPs, supporting message passing on factor graphs, policy inference
via expected free energy, and hierarchical generative models. It has
become the standard entry point for algorithm development.

\textbf{SPM.} The SPM package (Wellcome Centre for Human Neuroimaging)
includes MATLAB implementations of Dynamic Causal Modeling and
variational Bayesian inference under the FEP. It remains the reference
implementation for neuroimaging applications.

\textbf{RxInfer.jl.} RxInfer is a Julia package for reactive
message-passing-based Bayesian inference, supporting real-time and
streaming inference suitable for robotics and online learning. The
release of version 4.0.0 in early 2025 \citep{rxinfer2025} substantially
enhanced its probabilistic programming framework, introducing projected
constraints and adaptive qualities specifically optimized for dynamic
streams of data and autonomous systems.

\begin{block}{Comparative Feature Matrix}
\protect\phantomsection\label{comparative-feature-matrix}
{\def\LTcaptype{none} % do not increment counter
\begin{longtable}[]{@{}
  >{\raggedright\arraybackslash}p{(\linewidth - 6\tabcolsep) * \real{0.2500}}
  >{\raggedright\arraybackslash}p{(\linewidth - 6\tabcolsep) * \real{0.2500}}
  >{\raggedright\arraybackslash}p{(\linewidth - 6\tabcolsep) * \real{0.2500}}
  >{\raggedright\arraybackslash}p{(\linewidth - 6\tabcolsep) * \real{0.2500}}@{}}
\toprule\noalign{}
\begin{minipage}[b]{\linewidth}\raggedright
Feature
\end{minipage} & \begin{minipage}[b]{\linewidth}\raggedright
pymdp
\end{minipage} & \begin{minipage}[b]{\linewidth}\raggedright
SPM
\end{minipage} & \begin{minipage}[b]{\linewidth}\raggedright
RxInfer.jl
\end{minipage} \\
\midrule\noalign{}
\endhead
\textbf{Language} & Python & MATLAB & Julia \\
\textbf{State Spaces} & Discrete & Discrete + Continuous & Continuous
(factor graphs) \\
\textbf{Inference} & Message passing & Variational Bayes & Reactive
message passing \\
\textbf{Deep AIF} & Partial & No & Via custom factors \\
\textbf{Real-time} & No & No & Yes (streaming) \\
\textbf{Hierarchical} & Yes & Yes (DCM) & Yes \\
\textbf{License} & MIT & GPL & MIT \\
\textbf{Primary Use} & Research prototyping & Neuroimaging & Robotics /
online learning \\
\bottomrule\noalign{}
\end{longtable}
}

The complementary strengths across these packages reveal a structurally
fragmented ecosystem: \texttt{pymdp} provides an accessible,
Python-native entry point for discrete prototyping; \texttt{SPM} remains
the clinical gold standard for continuous neuroimaging; and
\texttt{RxInfer.jl} addresses the real-time constraints of embedded
robotics. The absence of a unified, cross-regime computational
infrastructure represents both a critical operational bottleneck and a
major opportunity for framework unification.
\end{block}
\end{block}

\begin{block}{Knowledge Graph Infrastructure}
\protect\phantomsection\label{knowledge-graph-infrastructure}
Our knowledge graph uses an RDF-compatible schema deployable on standard
semantic web infrastructure. The engineering trade-offs among the three
deployment options are straightforward:

\textbf{Nanopublication servers} provide decentralized,
content-addressed storage. The pipeline writes nanopublications in two
forms: JSON Lines (for incremental checkpointing and tooling) and
RDF/TriG per the \href{https://nanopub.net/}{nanopublication standard}
(Assertion, Provenance, Publication Info), suitable for the
nanopublication network and FAIR deployment.

\textbf{RDF stores} (e.g., Apache Jena Fuseki, Blazegraph, Oxigraph)
enable SPARQL queries such as ``find all papers supporting hypothesis H
published after 2020 in the neuroscience domain (C1).'' The cost is
operational overhead and query latency.

\textbf{Property graph databases} (e.g., Neo4j) prioritize traversal
performance for path queries and community detection, at the expense of
semantic web compatibility.

The namespace \texttt{http://activeinference.org/ontology/} ensures
integration with external ontologies and linked data resources.
\end{block}

\begin{block}{Multi-Level Quality Assurance}
\protect\phantomsection\label{multi-level-quality-assurance}
Quality assurance operates at four levels.

\begin{block}{Assertion-Level Validation}
\protect\phantomsection\label{assertion-level-validation}
Assertions below a configurable confidence threshold (default 0.5) are
flagged for review. Inter-annotator agreement (\(\kappa\)) is computed
when multiple annotators assess the same paper.
\end{block}

\begin{block}{Graph-Level Consistency Checks}
\protect\phantomsection\label{graph-level-consistency-checks}
Consistency checks verify that all nodes link to valid targets and no
orphan nodes exist. Coverage metrics track the proportion of annotated
papers.
\end{block}

\begin{block}{Score-Level Unit Testing}
\protect\phantomsection\label{score-level-unit-testing}
Hypothesis scoring is validated through unit tests with synthetic data
verifying boundary conditions (all-support → +1, all-contradict → −1,
balanced → 0). Sensitivity analysis varies confidence thresholds and
citation weighting.
\end{block}

\begin{block}{Pipeline-Level Test Coverage}
\protect\phantomsection\label{pipeline-level-test-coverage}
Test-driven development enforces 90\% minimum code coverage on project
modules and 60\% on shared infrastructure, with real data and
computation (no mocking).
\end{block}

\begin{block}{Quality Thresholds}
\protect\phantomsection\label{quality-thresholds}
{\def\LTcaptype{none} % do not increment counter
\begin{longtable}[]{@{}llll@{}}
\toprule\noalign{}
Level & Metric & Threshold & On Failure \\
\midrule\noalign{}
\endhead
Assertion & Confidence & \(\geq 0.5\) & Flag for review \\
Assertion & Inter-annotator \(\kappa\) & \(\geq 0.6\) & Re-annotate \\
Graph & Orphan node ratio & \(= 0\) & Reject build \\
Graph & Corpus coverage & \(\geq 80\%\) & Warning \\
Score & Boundary tests & All pass & Block release \\
Pipeline & Code coverage & \(\geq 90\%\) & Block merge \\
Pipeline & Test pass rate & \(100\%\) & Block release \\
\bottomrule\noalign{}
\end{longtable}
}

The hypothesis evidence results, temporal dynamics of evidence
accumulation, and assertion analysis are presented in the dedicated
hypothesis results section (see
\hyperref[sec:hypothesis_results]{Section 6}).
\end{block}
\end{block}
\end{frame}

\end{document}
